% =========================================================
% capitulos/02_marco_teorico.tex
% Placeholder con las herramientas que se van a usar en la
% tesis: tablas con el estilo azul, pseudocódigo y ecuaciones
% numeradas / sin numerar de distintos tipos.
% =========================================================

\chapter{Marco Teórico}
\label{chap:marco_teorico}


%%%% AISLAR %%%%%

\section{Representación y visualización de datos multidimensionales}
\label{sec:Representacion_Visualizacion}

En el desarrollo de sistemas de aprendizaje automático, el análisis exploratorio de los datos es una etapa fundamental para comprender la naturaleza del problema. Antes de entrenar modelos de clasificación complejos, la capacidad de observar geométricamente cómo se distribuyen las observaciones permite explorar patrones, detectar posibles valores atípicos e inspeccionar el grado de solapamiento entre las distintas clases \cite{oettl2024artificial,abukmeil2021survey}. 

En este contexto resulta conveniente distinguir entre la visualización y 
la representación que la hace posible. La visualización constituye la manifestación gráfica de los datos, mientras que la representación corresponde al espacio de dimensión reducida utilizado para proyectar las observaciones y facilitar su inspección.



\subsection{El desafío de visualizar datos de alta dimensión}

La representación visual directa de datos multidimensionales está limitada por la capacidad de representar simultáneamente un número reducido de dimensiones, por lo que las representaciones visuales convencionales suelen restringirse a una, dos o tres dimensiones \cite{andrews1972plots,espadoto2021unprojection} .

% CORREGIR CITA ARRIBA
 

Intentar comprender la estructura de un conjunto de datos multidimensional a través del análisis aislado por pares de variables resulta ineficiente y propenso a omitir interacciones complejas que sean clave para entenderlos \cite{neto2022multivariate}. A medida que la dimensionalidad aumenta, los datos se vuelven escasos en el espacio representativo, agravando los efectos de la maldición de la dimensionalidad \cite{peng2023interpreting}. Por lo tanto, surge la necesidad de proyectar esta información hacia espacios de menor dimensión que permitan conservar, de acuerdo con un criterio determinado, características relevantes de la estructura de los datos.


\subsection{Reducción dimensional como mecanismo de representación}

La reducción dimensional comprende un conjunto de técnicas orientadas a representar los datos en un espacio de baja dimensión, con el propósito de preservar, conforme al criterio de la técnica empleada, aquellas propiedades consideradas relevantes de la estructura original \cite{zebari2020comprehensive}. 

Cuando la dimensión resultante es suficientemente baja, estas representaciones también pueden emplearse para la exploración visual de los datos \cite{hasan2021review}.

En esta tesis se consideran representaciones de baja dimensión aquellas cuya dimensión no sea mayor a tres, es decir, $d\leq3$.


Este proceso de representación puede lograrse seleccionando un subconjunto específico de variables (selección de características) o transformando el espacio mediante combinaciones matemáticas (extracción de características) \cite{nogales2023analysis}. 

La selección de características conserva un subconjunto de variables originales mientras que la extracción construye nuevas dimensiones mediante transformaciones o combinaciones de las variables disponibles. En el contexto de la visualización, ambas estrategias pueden producir representaciones de baja dimensión, aunque difieren en el grado de correspondencia directa que mantienen con las variables originales \cite{2}.

\subsection{Proyecciones discriminantes para la visualización supervisada}

La utilidad de una visualización depende fuertemente del criterio matemático utilizado para proyectar los datos. El Análisis de Componentes Principales (PCA) constituye una técnica ampliamente utilizada para obtener representaciones de baja dimensión. Al proyectar las observaciones sobre sus primeras componentes, se obtiene un mapa que preserva la mayor varianza global de los datos. Sin embargo, al ser un método no supervisado, PCA no incorpora las etiquetas de clase en el proceso de obtención de las componentes \cite{PENDIENTE_LimitacionesPCAVisualizacion}.


En problemas de clasificación, la varianza global no siempre coincide con la separabilidad de las clases. Cuando el objetivo específico es favorecer la separación entre grupos, resulta conveniente emplear proyecciones discriminantes supervisadas, cuyo propósito es orientar la proyección hacia direcciones que favorezcan la separación entre los grupos \cite{CORREGIR}.

Para ilustrar esta diferencia fundamental entre la maximización de la varianza global y la maximización de la separabilidad, la figura \ref{fig:MT1} presenta una comparativa geométrica conceptual. Mientras que PCA proyecta los datos buscando la dirección de mayor dispersión general (ignorando las etiquetas de clase), una proyección discriminante orienta los ejes específicamente para maximizar la separación entre las clases respecto a su dispersión interna \cite{}.

\begin{figure}[htbp]
	\centering
	\includegraphics[width=1\textwidth]{MT1.png}
	\caption{Comparativa conceptual entre los criterios de proyección de PCA y LDA. Elaboración propia.}
	\label{fig:MT1}	
\end{figure}

\subsection{Visualización de representaciones de baja dimensión}


Las representaciones de baja dimensión permiten convertir un conjunto de observaciones definido originalmente en un espacio multidimensional en una geometría que puede ser inspeccionada visualmente. Cuando la representación tiene una, dos o tres dimensiones, es posible recurrir a recursos gráficos que permiten explorar directamente la distribución de las observaciones y su relación con las clases \cite{}.

En un espacio proyectado de baja dimensión, una separación visualmente marcada entre los grupos puede ofrecer evidencia exploratoria acerca de la existencia de una estructura discriminante; sin embargo, dicha representación no constituye por sí misma una prueba estadística de separabilidad ni garantiza desempeño predictivo \cite{}. Su valor radica en proporcionar una representación geométrica que facilite la exploración de la estructura de los datos y la formulación de hipótesis sobre la relación entre las clases, el grado de solapamiento y la presencia de observaciones atípicas, mientras que la evaluación de la capacidad discriminante y del rendimiento predictivo de un clasificador debe sustentarse en métricas cuantitativas \cite{}.

Además de la posibilidad de inspeccionar visualmente la estructura proyectada, una representación de baja dimensión puede adquirir mayor utilidad interpretativa cuando es posible establecer una relación explícita entre sus dimensiones y las características que intervienen en su construcción \cite{}.

Desde esta perspectiva, la trazabilidad de los ejes resulta relevante para determinar que características participan en la construcción de la representación. 
En ciertos métodos de reducción dimensional, especialmente aquellos que recurren a transformaciones implícitas del espacio de características, la relación explícita entre los ejes obtenidos y las variables originales tiende a ser menos directa, lo que dificulta la interpretación de la representación en términos de las características iniciales \cite{}.



La necesidad de contar con representaciones de baja dimensión que, además de facilitar la inspección visual, incorporen explícitamente la información sobre la pertenencia de las observaciones a distintas clases, motiva el empleo de técnicas de proyección discriminante. En este contexto, LDA ofrece un marco matemático que permite construir ejes orientados a maximizar la separación entre clases, estableciendo así la base de una representación discriminante supervisada \cite{}.





%%%% AISLAR %%%%%

\section{Análisis Discriminante Lineal}

\subsection{Fundamentos y criterio de Fisher}

El Análisis Discriminante Lineal, comúnmente referido como LDA por sus siglas en inglés, es una técnica estadística multivariante propuesta originalmente por Fisher (1936), cuyo propósito es encontrar combinaciones lineales de las variables originales que permitan favorecer la separación entre dos o más clases. En el contexto de la reducción de dimensionalidad supervisada, LDA busca proyectar los datos desde un espacio original de dimensión ($p$) hacia un subespacio de menor dimensión ($d$), procurando que las observaciones pertenecientes a una misma clase permanezcan compactas, mientras que las clases distintas se encuentren lo más separadas posible \cite{CORREGIR}.

A diferencia de PCA, LDA incorpora explícitamente la información de clase en la construcción de la proyección, orientando sus direcciones hacia la separación entre los grupos  \cite{bishop2006pattern}.

Desde una perspectiva probabilística, LDA suele asociarse con el supuesto de que las densidades condicionales de las clases siguen distribuciones normales multivariantes con una matriz de covarianza común, es decir, bajo una condición de homocedasticidad \cite{mclachlan1992discriminant}. Sin embargo, desde la formulación geométrica de Fisher, el método también puede entenderse como la búsqueda de una proyección lineal que maximiza una razón de separabilidad entre clases \cite{fisher1936use}. Esta distinción es importante, ya que en aplicaciones prácticas LDA puede emplearse tanto como clasificador como técnica de reducción de dimensionalidad supervisada \cite{martinez2001pca}.

Para introducir la formulación básica, considérese primero un problema de clasificación binaria con dos clases $C_1$ y $C_2$. El objetivo consiste en encontrar un vector de proyección $w$ que transforme cada observación del espacio original en un valor escalar mediante una combinación lineal de sus características. La orientación óptima de este vector se determina buscando que las medias de las clases proyectadas queden suficientemente separadas, al mismo tiempo que se controla la variabilidad interna de cada clase.

\begin{equation}
	y = \mathbf{w}^Tx
	\label{eq:EMT 1}
\end{equation}

Para formalizar esta idea, se definen dos cantidades principales: la matriz de dispersión intraclase ($S_W$) y la matriz de dispersión inter-clase ($S_B$). La primera resume la variabilidad de las observaciones alrededor de sus respectivos centroides de clase, mientras que la segunda cuantifica la separación entre los centroides de las clases. 

\begin{equation}
	\mathbf{S}_W = \sum_{n \in C_1} (\mathbf{x}_n - \mathbf{m}_1)(\mathbf{x}_n - \mathbf{m}_1)^T + \sum_{n \in C_2} (\mathbf{x}_n - \mathbf{m}_2)(\mathbf{x}_n - \mathbf{m}_2)^T
	\label{eq:EMT2}
\end{equation}

A partir de estas matrices, el criterio de Fisher se expresa como una razón entre la dispersión interclase y la dispersión intraclase, de modo que la dirección discriminante óptima es aquella que maximiza dicha razón.


\begin{equation}
	J\mathbf{(w)} = \mathbf{\frac{w^TS_Bw}{w^TS_Ww}}	
	\label{eq:EMT3}
\end{equation}

En el caso binario, la solución analítica conduce a una dirección de proyección proporcional al producto entre la inversa de la matriz de dispersión intraclase y la diferencia entre los vectores de medias de ambas clases. Esta formulación permite entender la intuición central del método: proyectar los datos hacia una dirección donde las clases se separen tanto como sea posible, manteniendo simultáneamente baja la dispersión interna de cada grupo \cite{bishop2006pattern}.

\begin{equation}
	\mathbf{w^*} \propto \mathbf{S_W^{-1}(m_1-m_2)}
	\label{eq:EMT4}
\end{equation}

\subsection{Extensión multiclase y obtención de ejes discriminantes}

Aunque la formulación binaria permite presentar la intuición geométrica de LDA, muchos problemas de clasificación involucran tres o más clases. En el caso multiclase, el objetivo general se mantiene: encontrar un conjunto de direcciones discriminantes que maximicen la separación entre clases y minimicen la dispersión dentro de cada clase. Sin embargo, en lugar de obtener una única dirección de proyección, se busca una matriz de proyección formada por varios vectores discriminantes.

Para ello, se define primero la media global del conjunto de datos, así como los vectores de medias correspondientes a cada clase. Con base en estos centroides, la matriz de dispersión intraclase acumula la variabilidad interna de cada grupo, mientras que la matriz de dispersión interclase acumula la separación ponderada entre cada centroide de clase y la media global del conjunto de datos.

\begin{equation}
	\mathbf{m} = \frac{1}{N} \sum_{k=1}^K N_k \mathbf{m}_k, \qquad \mathbf{S}_W = \sum_{k=1}^K \mathbf{S}_k, \qquad \mathbf{S}_B = \sum_{k=1}^K N_k (\mathbf{m}_k - \mathbf{m})(\mathbf{m}_k - \mathbf{m})^T
	\label{eq:EMT5}
\end{equation}

En esta formulación, el criterio de optimización busca maximizar la razón entre la dispersión interclase y la dispersión intraclase en el espacio proyectado. Una forma común de expresar este criterio es mediante la razón entre los determinantes de las matrices de dispersión proyectadas \cite{fukunaga2013introduction}.

\begin{equation}
	J(\mathbf{W}) = \frac{|\mathbf{W}^T \mathbf{S}_B \mathbf{W}|}{|\mathbf{W}^T \mathbf{S}_W \mathbf{W}|}
	\label{eq:EMT6}
\end{equation}

La maximización de este criterio conduce a un problema generalizado de valores propios, en el cual los vectores propios definen las direcciones discriminantes y los valores propios asociados indican la relevancia discriminante de cada dirección \cite{martinez2001pca}.

\begin{equation}
	\mathbf{S_Bw_i = \lambda_iS_Ww_i}
	\label{eq:EMT7}
\end{equation}

Los vectores propios asociados a los valores propios más grandes se seleccionan para formar la matriz de proyección final. En este sentido, cada valor propio puede interpretarse como una medida de la separabilidad capturada por su dirección discriminante correspondiente.

Por ello, LDA no maximiza “los ejes” en sí mismos, sino un criterio de separabilidad cuya solución produce los ejes discriminantes del nuevo espacio.


Un aspecto importante de la formulación multiclase es que el número máximo de dimensiones discriminantes que pueden extraerse está limitado por el número de clases. En particular, si existen ($C$) clases, el número máximo de direcciones discriminantes útiles es ($C-1$) \cite{duda2000pattern}. Esto se debe a que la matriz de dispersión interclase tiene un rango máximo determinado por la cantidad de centroides de clase disponibles. Por tanto, LDA no solo reduce la dimensionalidad, sino que también impone una estructura supervisada al espacio reducido \cite{belhumeur1997eigenfaces}.

\subsection{Supuestos, limitaciones y comparación con QDA}

La formulación clásica de LDA descansa sobre varios supuestos que condicionan su comportamiento. Desde el punto de vista probabilístico, el método asume que las clases pueden modelarse mediante distribuciones normales multivariantes con una matriz de covarianza común. Esta hipótesis conduce a fronteras de decisión lineales. Desde el punto de vista geométrico, LDA asume que una proyección lineal es suficiente para separar adecuadamente las clases en el espacio transformado.

Estas condiciones explican algunas de sus principales limitaciones. Cuando las clases presentan fronteras de decisión no lineales, distribuciones multimodales o patrones de separación complejos, una frontera lineal puede resultar insuficiente. En tales casos, la proyección obtenida por LDA puede no reflejar adecuadamente la separabilidad real de los datos, lo que limita su desempeño como clasificador o como técnica de reducción de dimensionalidad supervisada \cite{mika1999fisher}.

Otra limitación importante se relaciona con la necesidad de invertir o resolver sistemas asociados a la matriz de dispersión intraclase. Cuando el número de variables es comparable o superior al número de observaciones, la matriz de dispersión intraclase puede volverse singular o mal condicionada. Este escenario, conocido como problema de muestras pequeñas (Small Sample Size problem), impide aplicar directamente la formulación clásica de LDA y puede producir inestabilidad numérica en la obtención de los vectores discriminantes \cite{gao2006direct}.

Una alternativa clásica frente a la rigidez lineal de LDA es el Análisis Discriminante Cuadrático (QDA). A diferencia de LDA, que estima una matriz de covarianza común para todas las clases, QDA estima una matriz de covarianza distinta para cada clase. Esta diferencia permite construir fronteras de decisión cuadráticas, lo que puede resultar útil cuando las clases presentan estructuras de dispersión diferentes. Sin embargo, QDA debe entenderse principalmente como un clasificador probabilístico y no como un método de proyección supervisada de baja dimensión. Además, al requerir la estimación de una matriz de covarianza por clase, puede ser más sensible que LDA en escenarios de alta dimensionalidad y tamaño muestral reducido \cite{duda2000pattern}.

\begin{equation}
	\delta_k(x) = -\frac{1}{2} \ln |\Sigma_k| - \frac{1}{2} (x - \mu_k)^T \Sigma_k^{-1} (x - \mu_k) + \ln \pi_k
	\label{eq:EMT8}
\end{equation}

En consecuencia, aunque QDA ofrece mayor flexibilidad para modelar fronteras de decisión no lineales, no resuelve por sí mismo el problema de reducción de dimensionalidad supervisada ni elimina las dificultades asociadas a la estimación de matrices de covarianza en escenarios ($p>n$). Por ello, su inclusión resulta pertinente como método comparativo, pero no sustituye la necesidad de estrategias orientadas a estabilizar o reformular la proyección discriminante.

\subsection{Implementación computacional y tratamiento de la singularidad}

La formulación algebraica clásica de LDA suele expresarse a partir de la inversión de la matriz de dispersión intraclase ($\mathbf{S_W}$) o, de forma equivalente, mediante un problema generalizado de valores propios. Esta formulación resulta problemática cuando ($\mathbf{S_W}$) es singular o está mal condicionada, situación frecuente en escenarios de alta dimensionalidad, presencia de multicolinealidad o tamaño muestral reducido. En estos casos, el cálculo directo de ($\mathbf{S_W^{-1}}$) puede ser numéricamente inestable o incluso inviable \cite{friedman1989regularized}.

En la práctica, algunas implementaciones computacionales de LDA reformulan el problema para evitar la inversión explícita de la matriz de dispersión intraclase. 

Este es el caso de la función lda del paquete MASS en R, cuya implementación construye una transformación de escala que lleva las observaciones hacia funciones discriminantes normalizadas de manera que la covarianza intragrupo sea esférica.

Esta operación puede interpretarse como una forma de blanqueamiento del espacio, ya que busca remover la estructura de correlación interna de las clases antes de estimar las direcciones discriminantes \cite{venables2002modern}.

Matemáticamente, el proceso de blanqueamiento o esferización busca encontrar una matriz de transformación $\mathbf{W}$ tal que, al proyectar el espacio original, la nueva matriz de dispersión intraclase se convierta en la matriz identidad:

\begin{equation}
	\mathbf{W^TS_WW = I}
	\label{eq:EMT9}
\end{equation}

Mediante la descomposición espectral se tiene que 

\begin{equation}
	\mathbf{S}_W = \mathbf{U} \Lambda \mathbf{U}^T,
	\label{eq:EMT10}
\end{equation}

donde $\mathbf{U}$ es la matriz ortogonal de vectores propios y $\Lambda$ es la matriz diagonal de valores propios asociados. La matriz transformadora se define entonces como $\mathbf{W} = \mathbf{U} \Lambda^{-1/2}$.

Al proyectar los datos originales empleando $\mathbf{W}$, la varianza interna de los grupos se vuelve uniforme en todas las direcciones (esférica), reduciendo la estructura de correlación original y estabilizando los cálculos posteriores.

Cuando $\mathbf{S_W}$ es singular, esta transformación debe entenderse sobre el subespacio asociado a valores propios positivos o mediante una versión truncada/pseudoinversa, ya que los valores propios nulos no pueden invertirse.

La implementación utiliza descomposición en valores singulares (SVD) para identificar las direcciones discriminantes relevantes. En lugar de depender de la inversión directa de ($\mathbf{S_W}$), la SVD permite trabajar con el rango efectivo de la matriz y descartar direcciones asociadas con varianza nula o numéricamente despreciable. De esta forma, cuando existen variables colineales o combinaciones lineales con varianza muy baja, el algoritmo puede detectarlas mediante un umbral de tolerancia y restringir el cálculo a las dimensiones que conservan información discriminativa suficiente.

Formalmente, la Descomposición en Valores Singulares factoriza una matriz de datos $\mathbf{X}$ (previamente centrada y blanqueada) de la forma $\mathbf{X} = U \Sigma V^T$, donde $\Sigma$ es una matriz diagonal que contiene los valores singulares $\sigma_i$. 

El rango efectivo de la matriz se determina mediante el análisis de estos valores singulares, descartando aquellos que caen por debajo de un umbral de tolerancia $\tau$ cercano a cero ($\sigma_i < \tau$). Al truncar la descomposición y conservar únicamente las dimensiones con varianza significativa, el algoritmo evita la división por cero o la inestabilidad numérica que provocarían las dimensiones colineales al calcular las proyecciones discriminantes finales \cite{golub2013matrix}.

Esta estrategia no elimina por completo el problema Small Sample Size, ya que la disponibilidad de direcciones discriminantes sigue dependiendo de la información contenida en los datos y del rango efectivo de las matrices involucradas. Sin embargo, puede obtener una proyección en muchos escenarios donde la formulación clásica sería inestable, siempre que exista rango efectivo suficiente \cite{krzanowski1995discriminant}.

En el contexto de esta investigación, esta consideración es relevante porque la expansión polinomial incrementa el número de variables y puede introducir dependencias lineales o casi lineales entre términos originales, términos cuadrados e interacciones \cite{montgomery2022introduction}. 

Por ello, utilizar una implementación de LDA basada en blanqueamiento, control de singularidad y SVD permite emplear el discriminante lineal como motor de evaluación dentro de un esquema wrapper, siempre que previamente se controle la dimensionalidad mediante selección evolutiva. Así, la estabilidad computacional de LDA no se asume como una solución aislada al problema SSS, sino como un componente operativo que se beneficia de la reducción previa del espacio de características.

\section{Problema de muestras pequeñas (Small Sample Size)}

El de muestras pequeñas o Small Sample Size (SSS) aparece en escenarios donde el número de variables descriptivas es elevado en relación con el número de observaciones disponibles. En términos generales, esta situación suele expresarse como ($p>n$), donde ($p$) representa la dimensión del espacio de características y ($n$) el número de muestras. Este escenario es frecuente en dominios de alta dimensionalidad, como bioinformática, análisis de señales, reconocimiento de patrones y problemas donde se generan nuevas variables mediante transformaciones del espacio original \cite{raudys1991small}\cite{sorzano2014survey}.

En métodos basados en la estimación de matrices de covarianza o dispersión, el problema SSS se manifiesta matemáticamente como una deficiencia de rango. En el caso del Análisis Discriminante Lineal (LDA), la matriz de dispersión intraclase ($\mathbf{S_W}$) tiene dimensión ($p\times p$), pero su rango máximo está limitado por la cantidad de observaciones disponibles y por el número de clases. En particular, para ($C$) clases, se cumple que:

\begin{equation}
	\mathbf{rank(S_W)}\leq n-C
	\label{eq:EMT11}
\end{equation}

Por tanto, bajo la condición de singularidad donde $p>n-C$, la matriz $\mathbf{S_W}$ no puede alcanzar un rango completo. Esta deficiencia la vuelve singular (no invertible), lo que impide aplicar directamente la formulación clásica que requiere el cálculo algebraico de $\mathbf{S_W^{-1}}$.

La consecuencia práctica de este problema es que la estimación de las direcciones discriminantes puede volverse numéricamente inestable o no estar definida bajo la formulación algebraica clásica. En LDA, esto afecta la resolución del problema generalizado de valores propios y limita la obtención de proyecciones discriminantes estables \cite{gao2006direct} \cite{martinez2001pca}. Aunque existen estrategias computacionales para reducir la dependencia de la inversión directa de $\mathbf{S_W}$, como el uso de la Descomposición en Valores Singulares (SVD), el blanqueamiento o la regularización, estas alternativas no eliminan por completo la dificultad intrínseca de trabajar con un espacio de variables que supera a la información disponible.

\section{Expansión polinomial y espacios de caracteristicas}

Ante las limitaciones geométricas del LDA descritas en la sección anterior, resulta pertinente transformar el espacio de representación de los datos. Esta investigación adopta el enfoque de Expansión de Funciones Base (Basis Function Expansion), una técnica ampliamente utilizada en el aprendizaje automático que permite extender la capacidad de los modelos lineales para representar relaciones no lineales \cite{bishop2006pattern}.

\subsection{Formalización de la expansión polinomial}

El procedimiento consiste en mapear el vector de entrada original $\mathbf{x} \in \mathbb{R}^d$ a un espacio de características de mayor dimensión $\mathcal{F} \in \mathbb{R}^D$ mediante una función no lineal $\phi(x)$.

Específicamente, se implementa una expansión polinomial de segundo orden (cuadrática), tal como se detalla en la literatura especializada \cite{kanade2016} \cite{niu2021local}.

Para un caso base bidimensional donde el vector de entrada es $x = [x_1,x_2]^T$, la función de mapeo $\phi(x)$ construye un nuevo vector que incluye el sesgo, los términos originales, los términos cuadráticos y las interacciones cruzadas:

\begin{equation}
	\phi(x) = [1,x_1,x_2,x_1^2,x_2^2,x_1x_2]^T
	\label{eq:EMT12}
\end{equation}

Para efectos prácticos de su uso en este trabajo, se omitirá el sesgo debido a que no aporta en el modelado de las fronteras no lineales.

\begin{figure}[htbp]
	\centering
	\includegraphics[width=0.6\textwidth]{MC2.png}
	\caption{Ejemplo de términos generados mediante expansión polinomial para dos variables.Elaboración propia.}
	\label{fig:MC2}
\end{figure}

Como se ilustra en la Figura \ref{fig:MC2}, cada componente cumple un rol discriminante específico:

\begin{itemize}
	\item Términos Lineales ($A,B$): preservan la información original.
	
	\item - Términos Cuadráticos ($A^2,B^2$): permiten incorporar curvatura en la
	representación y contribuyen a modelar fronteras de decisión no lineales.
	
	\item - Términos de Interacción ($AB$): representan efectos multiplicativos entre variables, permitiendo modelar relaciones que no pueden describirse mediante términos lineales independientes \cite{kanade2016}.
	
\end{itemize}

La justificación teórica para realizar esta expansión se fundamenta en el Teorema de Cover, el cual postula que un problema de clasificación complejo y no lineal en un espacio de baja dimensión puede aumentar la probabilidad de separabilidad lineal bajo una transformación no lineal hacia un espacio de mayor dimensión \cite{cover1965geometrical}.

\subsection{Limitaciones: Crecimiento Combinatorio y Complejidad}

Si bien la expansión polinomial permite representar relaciones no lineales, introduce un desafío crítico conocido como la Maldición de la Dimensionalidad. A medida que aumenta la dimensión original $d$ o el grado del polinomio $P$, el número de términos generados $D$ crece de forma combinatoria.

Para una expansión polinomial de grado $P$, la dimensión del espacio resultante $D$ está dada por la combinatoria con repetición:

\begin{equation}
	D = \binom{d+P}{P}-1
	\label{eq:EMT13}
\end{equation}

Tal como se observa en la Figura \ref{fig:MC3}, este crecimiento es explosivo. Por ejemplo, expandir un vector modesto de $d = 10$ variables a un grado cuadrático ($P=2$) genera 65 términos sin incluir el sesgo.

\begin{figure}[htbp]
	\centering
	\includegraphics[width=0.6\textwidth]{MC3.png}
	\caption{Crecimiento del número de atributos bajo una expansión polinomial de grado 2 sin término de sesgo. Fuente: Elaboración propia.}
	\label{fig:MC3}
\end{figure}

Este fenómeno conlleva dos consecuencias considerables descritas por \cite{bishop2006pattern}:

\begin{itemize}
	\item Costo Computacional: El cálculo de la matriz de covarianza inversa $\mathbf{S_W^{-1}}$ requerida por el LDA se vuelve computacionalmente costoso $(O(D^3))$.
	\item Sobreajuste (Overfitting): Al tener un exceso de parámetros frente a un número limitado de muestras, el modelo tiende a memorizar el ruido de los datos en lugar de aprender la estructura subyacente.
\end{itemize}

En el contexto de LDA, este crecimiento dimensional es especialmente relevante porque la matriz de dispersión intraclase pasa a definirse sobre el espacio expandido, por lo que su dimensión depende de ($D$) y no solo de la dimensión original ($d$). En consecuencia, si ($D>n-C$), la matriz ($S_W$) no puede alcanzar rango completo, reproduciendo o agravando el problema SSS descrito en la sección anterior. Por esta razón, la expansión polinomial debe acompañarse de mecanismos de selección de características que permitan conservar únicamente los términos con mayor aporte discriminativo y reducir el riesgo de singularidad, colinealidad y sobreajuste.

\section{Algoritmos genéticos}

Los algoritmos genéticos constituyen una familia de métodos de búsqueda y optimización inspirados en mecanismos evolutivos como selección, recombinación y mutación. En términos generales, estos algoritmos trabajan con una población de soluciones candidatas que se modifican iterativamente con el propósito de mejorar una función objetivo o función de aptitud \cite{Goldberg1988}.

Su utilidad en problemas de selección de características radica en que permiten explorar espacios combinatorios amplios sin evaluar exhaustivamente todos los subconjuntos posibles de variables. En el contexto de esta investigación, los algoritmos genéticos resultan especialmente pertinentes debido al crecimiento dimensional producido por la expansión polinomial. Una vez generados términos originales, términos cuadráticos e interacciones, el número de predictores puede crecer de manera acelerada, lo que aumenta el riesgo de singularidad, colinealidad y sobreajuste. Por ello, se requiere un mecanismo de selección capaz de identificar subconjuntos informativos dentro del espacio expandido. Desde esta perspectiva, el algoritmo genético opera como una estrategia de búsqueda dentro de un esquema wrapper, ya que la calidad de cada subconjunto se evalúa mediante el desempeño de un modelo discriminante \cite{kohavi1997wrappers}.

\subsection{Representación cromosómica y función de aptitud}

En selección de características, una solución candidata puede representarse mediante un cromosoma binario. Cada posición del cromosoma corresponde a una característica del espacio disponible, y su valor indica si dicha característica se incluye o se excluye del subconjunto evaluado. Matemáticamente, el cromosoma se define como un vector binario $z = [z_1,z_2,\dots,z_D]$, donde cada gen $z_j \in \{1,0\}$ y $D$ representa la dimensión total del espacio de características.

A modo de ejemplo, un cromosoma de cinco dimensiones estructurado como $z = [1,0,1,0,1]$ codifica un subconjunto que retiene exclusivamente las características $\{x_1,x_3,x_5\}$, descartando el resto.

En el caso de una expansión polinomial, cada gen puede representar una variable original, un término cuadrático o una interacción entre variables. Por ejemplo, si el espacio expandido contiene términos como $x_1,x_2,x_1^2,x_2^2$ y $x_1x_2$, el cromosoma determina cuáles de estos términos serán conservados para entrenar el modelo discriminante. Esta representación permite que el algoritmo genético explore combinaciones de términos que podrían no ser evidentes mediante criterios univariados de selección.

La función de aptitud define el criterio mediante el cual se evalúa la calidad de cada cromosoma. En un esquema wrapper, esta función suele depender del desempeño del clasificador acoplado al subconjunto seleccionado. Para evitar soluciones excesivamente grandes y penalizar la complejidad del modelo, la función de aptitud puede formalmente como \cite{siedlecki1989note}:

\begin{equation}
	fitness^{z} = \text{Desempeño}^{(z)}-\lambda \frac{|z|_0}{D}
	\label{eq:EMT14}
\end{equation}

En esta expresión, $\text{Desempeño}^{(z)}$ representa una métrica de clasificación (como la Exactitud o el $F1-Score$) obtenida con el subconjunto activo, $|z|_0$ denota la norma $\ell_0$ que contabiliza el número de características seleccionadas (genes con valor 1), $D$ es la dimensión máxima del espacio expandido y $\lambda$ es un hiperparámetro que controla la magnitud de la penalización por complejidad. 

De esta manera, el algoritmo busca optimizar un criterio combinado que maximice la capacidad predictiva y reducir la dimensionalidad efectiva del subconjunto evaluado.

\subsection{Ciclo evolutivo y operadores genéticos}

El proceso evolutivo sigue un ciclo iterativo estructurado, el cual se ilustra en el diagrama de flujo de la figura \ref{fig:MT2}. El ciclo comienza con la inicialización de una población de cromosomas, generada comúnmente de forma aleatoria para asegurar una exploración inicial diversa. En cada generación, los cromosomas son sometidos a una fase de evaluación mediante la función de aptitud y, posteriormente, se aplican los operadores genéticos (selección, cruza y mutación) para producir nuevas soluciones candidatas que conformarán la siguiente generación \cite{Goldberg1988}.

\begin{figure}[H]
	\centering
	\includegraphics[width=0.5\textwidth]{MT2.png} % Ajustar nombre del archivo
	\caption{Diagrama de flujo del ciclo evolutivo de un algoritmo genético. Tomada de \cite{us_algoritmos_geneticos}.}
	\label{fig:MT2}
\end{figure}

El operador de selección favorece la reproducción de los cromosomas con mayor aptitud. Su propósito es incrementar la probabilidad de que las soluciones más prometedoras transmitan su información genética a la siguiente iteración. Existen diferentes estrategias algorítmicas, como el torneo, la ruleta o el elitismo; sin embargo, todas comparten el objetivo guiar la búsqueda hacia regiones del espacio con mejores propiedades discriminativas.

La cruza o recombinación combina segmentos de dos cromosomas padres para generar descendencia. En la selección de características, este operador permite entrelazar subconjuntos parciales de variables provenientes de soluciones distintas, promoviendo la combinación de términos polinomiales. Un ejemplo de cruza de un punto, asumiendo un punto de corte en la tercera posición, se ilustra de la siguiente forma:

\begin{itemize}
	\item Padre 1: $z_{P1} = [1,0,1|0,1]$
	\item Padre 2: $z_{P2} = [0,1,0|1,0]$
	\item Hijo resultante: $z_{H} = [1,0,1|1,0]$ 
\end{itemize}

La mutación introduce alteraciones estocásticas en genes individuales, invirtiendo sus valores lógicos (de 0 a 1, o viceversa). Este operador es importante para preservar la diversidad dentro de la población y evadir la convergencia prematura en óptimos locales.

Un ejemplo de mutación uniforme sobre el tercer gen de un cromosoma se representa como: 

\begin{itemize}
	\item Cromosoma original: $z = [1,0,\textbf{1},0,1]$
	\item Cromosoma mutado: $z' = [1,0,\textbf{0},0,1]$
\end{itemize}

En conjunto, estos operadores permiten navegar de forma iterativa el vasto espacio de combinaciones posibles. El algoritmo concluye su ejecución al alcanzar un límite de generaciones, al agotar un presupuesto computacional o cuando el cambio en la aptitud intergeneracional cae por debajo de un umbral predefinido.

\subsection{Algoritmos genéticos como estrategia wrapper}

La aplicación de algoritmos genéticos a la selección de características puede entenderse como una estrategia wrapper, dado que cada subconjunto candidato se evalúa a partir del desempeño del modelo de aprendizaje final. A diferencia de los métodos de tipo filtro, los enfoques wrapper incorporan la retroalimentación directa del clasificador y permiten considerar la interacción entre las características seleccionadas y el comportamiento del modelo \cite{kohavi1997wrappers} \cite{guyon2003introduction}. Esta propiedad resulta especialmente relevante cuando los predictores incluyen términos polinomiales, ya que la utilidad de una interacción ($x_i,x_j$) puede depender de su combinación con otros términos del espacio expandido.

En la arquitectura propuesta para esta investigación, el algoritmo genético actúa como mecanismo de exploración sobre el espacio formado por variables originales y términos polinomiales. El objetivo principal no es someter el espacio expandido completo al análisis discriminante, sino identificar un subconjunto que conserve capacidad de separación interclase y reduzca la dimensionalidad efectiva del problema. De esta forma, la búsqueda evolutiva contribuye a mitigar los efectos del crecimiento combinatorio y del problema Small Sample Size.
La pertinencia de acoplar un algoritmo genético con LDA radica en que la selección previa de características puede mejorar la estabilidad algebraica del modelo discriminante. Si el espacio polinomial completo se utiliza sin reducción, la matriz de dispersión intraclase puede volverse singular o mal condicionada. En contraste, la reducción previa mediante la metaheurística puede favorecer una representación más compacta y manejable para la estimación discriminante. Por lo tanto, el algoritmo genético no actúa como sustituto del análisis discriminante, sino como un mecanismo previo de selección y reducción que facilita su aplicación sobre un espacio polinomial controlado.

\section{Métodos kernel y clasificadores no lineales}

Los métodos kernel constituyen una familia de enfoques ampliamente utilizados para
extender modelos lineales hacia escenarios no lineales. Su idea central consiste en
transformar los datos originales hacia un espacio de características de mayor dimensión, donde relaciones que no son linealmente separables en el espacio original pueden representarse mediante fronteras lineales en el espacio transformado. Esta lógica se relaciona con la motivación subyacente de la expansión polinomial: aumentar la capacidad de representación del modelo mediante transformaciones del espacio de entrada \cite{scholkopf2018learning}.

\subsection{Truco Kernel y funciones kernel}

Sea $x \in \mathbb{R}^p$ un vector de entrada. Un método kernel asume la existencia de una función de mapeo $\phi(x)$, que proyecta las observaciones desde el espacio original hacia un espacio de características posiblemente de alta o incluso infinita dimensión. En principio, el producto interno en dicho espacio transformado se define como:

\begin{equation}
	K(x_i,x_j) = \phi(x_i)^T\phi(x_j)
	\label{eq:EMT15}
\end{equation}

La ventaja computacional fundamental de los métodos kernel consiste en que no es necesario calcular explícitamente las coordenadas proyectadas $\phi(x)$. En su lugar, se define una función kernel $K(x_i,x_j)$ que evalúa directamente el producto interno entre dos observaciones en el espacio transformado. Este procedimiento algorítmico se conoce como el kernel trick y permite trabajar con espacios de características de alta dimensionalidad sin incurrir en el costo de construirlos de forma explícita \cite{bishop2006pattern}.

Entre las funciones kernel más estandarizadas en la literatura se encuentran el kernel lineal, el kernel polinomial y el kernel de función de base radial (RBF). El kernel lineal equivale al producto interno ordinario, mientras que el kernel polinomial permite incorporar interacciones de orden superior de forma implícita. Por su parte, el kernel RBF mide la similitud local entre observaciones e induce fronteras de decisión altamente flexibles.

Matemáticamente, se expresan de la siguiente manera:

\noindent\textbf{Kernel lineal:}
\begin{equation}
	K(\mathbf{x}, \mathbf{z}) = \mathbf{x}^T \mathbf{z}
	\label{eq:EMT16}
\end{equation}

\noindent\textbf{Kernel polinomial:}
\begin{equation}
	K(\mathbf{x}, \mathbf{z}) = (\gamma \mathbf{x}^T \mathbf{z} + c)^r
	\label{eq:EMT17}
\end{equation}

\noindent\textbf{Kernel RBF:}
\begin{equation}
	K(\mathbf{x}, \mathbf{z}) = \exp(-\gamma \|\mathbf{x} - \mathbf{z}\|^2)
	\label{eq:EMT18}
\end{equation}


En el contexto de esta tesis, el kernel polinomial resulta particularmente relevante por su vinculación teórica con la expansión polinomial explícita. La diferencia principal radica en que, mientras la expansión polinomial construye explícitamente los términos cuadráticos y las interacciones cruzadas, el kernel polinomial opera sobre un espacio expandido de forma implícita. Esta distinción es crítica porque el enfoque propuesto en esta investigación busca conservar una relación directa rastreable con los términos generados, favoreciendo la trazabilidad analítica de las variables seleccionadas.

\subsection{KFDA y SVM como enfoques no lineales de referencia}

El Análisis Discriminante de Fisher con Kernel (KFDA) se concibe como una extensión no lineal del criterio de Fisher clásico. Mientras que LDA busca direcciones discriminantes en el espacio original mediante combinaciones lineales, KFDA traslada el problema de optimización al espacio de características inducido por el kernel. En dicho espacio, el objetivo se mantiene: maximizar la separación interclase y minimizar la dispersión intraclase; no obstante, la solución algebraica se reformula en términos de productos kernel y matrices
de Gram \cite{mika1999fisher}.

Una diferencia relevante entre KFDA y otros clasificadores kernel radica en su relación con la reducción de dimensionalidad supervisada. Al heredar la lógica del criterio de Fisher, KFDA puede formularse como un método orientado a obtener ejes o proyecciones discriminantes en el espacio de características inducido por el kernel. En este sentido, mantiene una relación conceptual más cercana con LDA, aunque opera mediante una formulación dual basada en productos kernel. Por el contrario, aunque SVM también puede utilizar funciones kernel para inducir fronteras no lineales, su objetivo principal no es construir una proyección discriminante de baja dimensión, sino estimar una frontera de decisión con margen máximo para clasificar nuevas observaciones.

De manera análoga, las Máquinas de Vectores de Soporte (SVM) emplean el kernel trick para construir fronteras de decisión no lineales. En su formulación clásica, SVM busca un hiperplano de separación que maximice el margen geométrico entre las clases, apoyándose en las observaciones más relevantes, denominadas vectores de soporte \cite{cortes1995support}. Al integrar funciones kernel, la función de decisión para una nueva observación $x$ se formula de la siguiente manera general:

\begin{equation}
	f(x) = \sum_{i = 1}^{n} {\alpha_iy_iK(x_i,x) + b}
	\label{eq:EMT19}
\end{equation}

Para el marco experimental de esta investigación, SVM con kernel polinomial y SVM con kernel RBF se consideran los clasificadores no lineales de referencia. Ambos enfoques, sin embargo, presentan una alta sensibilidad a la parametrización, dependiendo de la correcta sintonización de hiperparámetros estructurales como el grado del polinomio $r$, el coeficiente $\gamma$, o la constante de regularización ($C$) \cite{chapelle2002choosing}.

\subsection{Ventajas, limitaciones e interpretabilidad}

Los métodos kernel presentan una ventaja importante: permiten introducir capacidad de modelado no lineal sin construir explícitamente todas las variables transformadas. En lugar de generar de manera directa los términos del espacio de características, estos métodos operan mediante productos internos definidos por una función kernel. Esta propiedad reduce la necesidad de manipular explícitamente espacios expandidos de alta dimensión y permite modelar fronteras de decisión no lineales estructuras de separación complejas en el espacio original.

No obstante, esta formulación también introduce limitaciones. Al permanecer el espacio transformado de forma implícita, la interpretación directa de las variables originales y de sus interacciones puede volverse menos transparente. En consecuencia, aunque el modelo puede alcanzar buen desempeño predictivo, puede ser más difícil identificar qué variables o combinaciones de variables explican la separación entre clases. Además, el desempeño de los métodos kernel depende de la elección de la función kernel y de la adecuada sintonización de sus hiperparámetros, lo cual puede afectar la capacidad de generalización, especialmente
en escenarios de tamaño muestral reducido \cite{chapelle2002choosing} \cite{rudin2019stop}. En contraste, la expansión polinomial explícita permite observar directamente los términos generados, como variables originales, términos cuadrados e interacciones. 

Esta característica favorece la trazabilidad de la representación transformada y permite identificar qué términos son conservados después del proceso de selección. Sin embargo, dicha ventaja tiene como costo el crecimiento combinatorio del número de características. Por ello, la arquitectura propuesta en esta investigación combina expansión polinomial explícita con selección evolutiva, buscando conservar términos informativos y controlar la dimensionalidad antes de aplicar el análisis discriminante.

En síntesis, los métodos kernel y las SVM proporcionan un marco de referencia relevante para el modelado no lineal. Su inclusión en este Marco Teórico permite contrastar dos estrategias distintas: la no linealidad implícita inducida por funciones kernel y la no linealidad explícita mediante expansión polinomial y selección de características. Esta diferencia resulta central para ubicar la propuesta de esta tesis frente a enfoques kernel tradicionales.
